\chapter{Conventions and notation}
\label{ch:conventions}

\newcommand{\cvMac}[1]{\texttt{\textbackslash #1}}

This chapter fixes the symbols, the typographical rules and the index conventions
that the rest of the compendium and the paper share. Everything here is a decision,
not a result. The decisions were taken once, in the file \code{paper/notation.tex},
and every chapter of this compendium is written against that file. A reader who
wants to check the paper line by line should read this chapter first, because several
of the paper's choices differ from the habits of the diffusion literature, and the
differences are deliberate.

\section{How to read this compendium}

Every chapter is self-contained at the level of a chapter: it defines the objects it
uses, derives every formula it states (the scalar case in full, then the multivariate
statement), and works one small example by hand. Four typographical devices recur.

\begin{itemize}
\item A \term{term} in bold at its first use is a definition; the term is also entered
  in the index at the end of the document.
\item A paragraph opening with \emph{Where it is used} closes each chapter and names
  the paper sections, the experiment identifiers (E01 to E21) and the code paths in
  which the chapter's objects appear.
\item A paragraph opening with \emph{Why it matters here} explains, inside a
  derivation, why a particular hypothesis or constant is the one the paper needs.
\item A path set in typewriter, such as \code{src/markov\_score/theory.py}, is a file
  of the repository. Every such path was checked to exist when the chapter was written.
\end{itemize}

Theorems, propositions, lemmas and corollaries are numbered within chapters, as are
definitions, examples, assumptions and remarks. ``The paper'' always means the ICML
manuscript \emph{What Survives Diffusion Noise?}; its sections, theorems and figures
are referred to in words (``the paper, Section 3, Theorem 2''), never by cross-reference,
because the two documents are compiled separately. Experiments are referred to by
their registry identifier (E01, \dots, E21), which names an entry of
\code{experiments/registry/}. No measured number from the experiments is typed in this
compendium: those numbers live in generated macros of the paper, and the compendium
says ``see the paper, Figure 2'' instead. Analytic constants such as $\tfrac16$,
$\tfrac1{24}$ or $e^{-4t}$ are typed freely.

\section{The symbol table}

\Cref{tab:symbols} lists every macro of \code{paper/notation.tex} and every macro
added in the compendium's preamble, with its rendering, its meaning and the chapter
in which the object is defined from scratch. The macro names are given so that a
reader of the sources can search for them.

{\small
\setlength{\tabcolsep}{4pt}
\begin{longtable}{@{}>{\raggedright\arraybackslash}p{2.6cm}>{\raggedright\arraybackslash}p{2.0cm}>{\raggedright\arraybackslash}p{7.1cm}>{\raggedright\arraybackslash}p{3.3cm}@{}}
\caption{Every macro of the paper's notation file and of the compendium preamble.}
\label{tab:symbols}\\
\toprule
macro & renders as & meaning & defined in \\
\midrule
\endfirsthead
\toprule
macro & renders as & meaning & defined in \\
\midrule
\endhead
\multicolumn{4}{@{}l}{\emph{Diffusion}}\\
\cvMac{Pz} & $\Pz$ & the clean law of a sequence, a density on $\Reals^{\nsites}$ & Markov chains \\
\cvMac{Pt} & $\Pt$ & the noised law at diffusion time $t$, the law of $\vx_t$ & The OU channel and Tweedie \\
\cvMac{Dt} & $\Dt$ & $1-e^{-2t}$, the conditional variance of one noised coordinate & The OU channel and Tweedie \\
\cvMac{score} & $\score$ & the diffusion score $\score(\vx,t)=\nabla_{\vx}\log\Pt(\vx)$ & The OU channel and Tweedie \\
\cvMac{gwn} & $\gwn$ & standardised Gaussian white noise, $\langle\eta_i(t)\eta_j(t')\rangle=\delta_{ij}\delta(t-t')$ & The OU channel and Tweedie \\
\cvMac{tmin}, \cvMac{tmax} & $\tmin$, $\tmax$ & the ends of the noise schedule & The protocol \\
\cvMac{pmean\{$\cdot$\}} & $\pmean{\va}$ & posterior mean $\E[\va\mid\vx_t=\vx]$, in bracket notation & The OU channel and Tweedie \\
\cvMac{pmeansite\{i\}} & $\pmeansite{i}$ & posterior mean of one site & The OU channel and Tweedie \\
\addlinespace
\multicolumn{4}{@{}l}{\emph{Vectors and matrices}}\\
\cvMac{vect\{v\}}, \cvMac{mat\{A\}} & $\vect{v}$, $\mat{A}$ & bold lowercase for vectors, bold uppercase for matrices and tensors & this chapter \\
\cvMac{va} & $\va$ & the clean chain, a point of $\Reals^{\nsites}$; its sites are the scalars $a_i$ & Markov chains \\
\cvMac{vx} & $\vx$ & the noised chain, $\vx=\vx_t\in\Reals^{\nsites}$ & The OU channel and Tweedie \\
\cvMac{vz} & $\vz$ & a standard normal draw in $\Reals^{\nsites}$ & Gaussian integrals and Hermite \\
\cvMac{Sig} & $\Sig$ & a covariance matrix; $\Sig_0$ is the clean chain's, $\Sig_t=e^{-2t}\Sig_0+\Dt\Id$ the noised one's & Markov chains \\
\cvMac{Id} & $\Id$ & the identity matrix & this chapter \\
\addlinespace
\multicolumn{4}{@{}l}{\emph{Chain and graphical model}}\\
\cvMac{nsites} & $\nsites$ & chain length, equal to the ambient dimension & Markov chains \\
\cvMac{nseq} & $\nseq$ & number of observed sequences & Markov chains \\
\cvMac{ngrid} & $\ngrid$ & number of quadrature grid points & Grid quadrature \\
\cvMac{gridpt\{k\}} & $\gridpt{k}$ & the $k$-th grid point & Grid quadrature \\
\cvMac{gridh} & $\gridh$ & the grid spacing & Grid quadrature \\
\cvMac{halfwidth} & $\halfwidth$ & the grid half-width; the domain is $[-\halfwidth,\halfwidth]$ & Grid quadrature \\
\cvMac{kernel} & $\kernel$ & the Markov transition kernel $\kernel(a'\mid a)$ & Markov chains \\
\cvMac{kernelth} & $\kernelth$ & a transition kernel with parameters $\params$ & Maximum likelihood and EM \\
\cvMac{like} & $\like$ & the unary likelihood factor $\like(x_i\mid a_i)$ & Factor graphs and BP \\
\cvMac{lmsg}, \cvMac{rmsg} & $\lmsg$, $\rmsg$ & forward and backward messages & Factor graphs and BP \\
\cvMac{innov} & $\innov$ & the innovation density & Markov chains \\
\cvMac{corr} & $\corr$ & the autoregression coefficient & Markov chains \\
\cvMac{ivar} & $\ivar$ & the innovation variance, $\ivar=1-\corr^2$ in the paper & Markov chains \\
\cvMac{ncomp} & $\ncomp$ & number of mixture components of the fitted innovation & Maximum likelihood and EM \\
\addlinespace
\multicolumn{4}{@{}l}{\emph{Estimation}}\\
\cvMac{params} & $\params$ & a parameter vector & Fisher information \\
\cvMac{Qfun} & $\Qfun$ & the EM auxiliary function & Maximum likelihood and EM \\
\cvMac{Xis} & $\Xis$ & the expected transition statistic, a matrix on the grid & Maximum likelihood and EM \\
\cvMac{Lik} & $\Lik$ & the marginal log-likelihood & Maximum likelihood and EM \\
\cvMac{mmse}, \cvMac{lmmse} & $\mmse$, $\lmmse$ & the posterior mean and the linear MMSE estimator & Gaussian closure and LMMSE \\
\addlinespace
\multicolumn{4}{@{}l}{\emph{Operators}}\\
\cvMac{E}, \cvMac{Prob} & $\E$, $\Prob$ & expectation, probability & Moments and cumulants \\
\cvMac{Reals} & $\Reals$ & the real line & this chapter \\
\cvMac{Cov}, \cvMac{Var} & $\Cov$, $\Var$ & covariance, variance & Moments and cumulants \\
\cvMac{KL} & $\KL$ & Kullback--Leibler divergence & Maximum likelihood and EM \\
\cvMac{normal\{m\}\{v\}} & $\normal{m}{v}$ & the Gaussian law with mean $m$ and variance $v$ & Gaussian integrals and Hermite \\
\cvMac{Pgauss} & $\Pgauss$ & Gaussian moment projection: the Gaussian with the argument's mean and variance & Gaussian closure and LMMSE \\
\cvMac{dd} & $\dd$ & the differential in integrals, $\dd x$ & this chapter \\
\addlinespace
\multicolumn{4}{@{}l}{\emph{Pending-data guard and repository pointers}}\\
\cvMac{needsdata} & red bracket & with an argument naming the slot: marks a place in the paper still waiting for a frozen number; the build fails if any survives & this chapter \\
\cvMac{codepath}, \cvMac{codefile} & typewriter & a repository path & this chapter \\
\cvMac{coderef\{p\}\{f\}} & path $\rightarrow$ name & a path with a function or target inside it & this chapter \\
\addlinespace
\multicolumn{4}{@{}l}{\emph{Macros the paper defines in \code{paper/main.tex}}}\\
\cvMac{Frob\{T\}} & $\Frob{\mat{T}}$ & Frobenius norm, the root of the sum of squares over ordered indices & Moments and cumulants \\
\cvMac{Hen\{k\}} & $\Hen{k}$ & the $k$-th probabilists' Hermite polynomial & Gaussian integrals and Hermite \\
\cvMac{kfour} & $\kfour$ & the fourth-cumulant tensor of the clean law & Moments and cumulants \\
\cvMac{Sgauss} & $\Sgauss$ & the score of the covariance-matched Gaussian twin, $-\Sig_t^{-1}\vx$ & Gaussian closure and LMMSE \\
\cvMac{tstar} & $\tstar$ & the crossover time $\tfrac12\log\lmax(\Sig_0)$ & Markov chains; Attenuation theorems \\
\cvMac{lmax} & $\lmax$ & the largest eigenvalue of a symmetric matrix & Markov chains \\
\cvMac{Pgaussl} & $\Pgaussl$ & the Gaussian law with the clean chain's covariance (the twin) & Gaussian closure and LMMSE \\
\addlinespace
\multicolumn{4}{@{}l}{\emph{Compendium-only macros}}\\
\cvMac{tr}, \cvMac{diag} & $\tr$, $\diag$ & trace; diagonal matrix built from a vector & this chapter \\
\cvMac{argmax}, \cvMac{argmin} & $\argmax$, $\argmin$ & maximiser and minimiser & this chapter \\
\cvMac{term\{$\cdot$\}} & bold & a defined term, also indexed & this chapter \\
\cvMac{where\{$\cdot$\}}, \cvMac{why\{$\cdot$\}} & italic lead-in & the two closing paragraphs described above & this chapter \\
\cvMac{code\{$\cdot$\}} & typewriter & a repository path or identifier & this chapter \\
\bottomrule
\end{longtable}}

Three macros are typeset slightly differently in the paper and in this compendium,
and the reader should not be surprised by the difference. The paper writes the
fourth-cumulant tensor as $\kappa_4$ in light face, the compendium as $\kfour$ in
bold, because the compendium's rule is that tensors are bold. The paper writes the
twin's score with a superscript, $S^{\mathrm{G}}$, the compendium with a subscript,
$\Sgauss$. The paper writes the crossover as $t_\ast$ and the compendium as $\tstar$.
The objects are the same.

\section{Indices: sites, sequences, grid points, and why $n$ and $d$ are not used}

\term{Sites} are the coordinates of a sequence. A clean sequence is
$\va=(a_1,\dots,a_{\nsites})$, and $i$, $j$ always index sites, $i\in\{1,\dots,\nsites\}$.
The \term{chain length} $\nsites$ is therefore also the \term{ambient dimension} of
the space in which the diffusion runs: the noised sequence $\vx_t$ lives in
$\Reals^{\nsites}$, the score is a map $\Reals^{\nsites}\times[0,\infty)\to\Reals^{\nsites}$,
and every covariance matrix is $\nsites\times\nsites$. The paper writes this
identification once, as $n=\nsites$, and then never uses $n$ again.

\term{Sequences} are the independent draws. The data are $\nseq$ sequences
$\va^{1},\dots,\va^{\nseq}$ (or their noised versions), and $\mu$ indexes sequences,
$\mu\in\{1,\dots,\nseq\}$. Fisher information is stated \emph{per sequence}, and the
Cram\'er--Rao and local asymptotic minimax statements are in the limit
$\nseq\to\infty$ at fixed $\nsites$ and fixed $t$.

\term{Grid points} are the abscissae on which messages are stored. There are
$\ngrid$ of them, $\gridpt{1},\dots,\gridpt{\ngrid}$, uniformly spaced by $\gridh$ on
$[-\halfwidth,\halfwidth]$, and $k$, $l$ index them. A matrix on the grid, such as the
discretised kernel $\kernel(\gridpt{k}\mid\gridpt{l})$ or the expected transition
statistic $\Xis_{kl}$, is $\ngrid\times\ngrid$ and is indexed \emph{child first, parent
second}: $k$ is the site $a_i$, $l$ is the site $a_{i-1}$. This orientation is the one
thing a self-consistent recursion cannot catch if it is wrong, which is why
\code{tests/test\_fisher\_bruteforce.py} checks it against an enumeration that shares
nothing with the recursion.

\term{Multi-indices} $\nu=(\nu_1,\dots,\nu_{\nsites})\in\mathbb{N}^{\nsites}$ label
Hermite products; the \term{degree} is $|\nu|=\sum_i\nu_i$ and $\nu!=\prod_i\nu_i!$.

Why $n$ and $d$ are avoided. Much of the diffusion literature writes $n$ for the number
of samples and $d$ for the dimension of the data, and much of the statistics
literature writes $n$ for the sample size as well. In this laboratory the dimension
\emph{is} the chain length, and the reader should never have to remember that a
statement ``in dimension $d$'' is a statement about a chain of $d$ sites; so the
dimension is called $\nsites$ and the number of sequences $\nseq$. The letter $n$ is
kept for the degree of a scalar Hermite polynomial $\Hen{n}$ and for nothing else.
Two further collisions were removed at the same time: the transition kernel is $\kernel$
and only $\kernel$, and the number of grid points is $\ngrid$, whereas an earlier draft
used $K$ for the grid size and $W$ for the kernel.

Levels. The noise schedule is a finite set of diffusion times, twelve of them
log-spaced on $[\tmin,\tmax]$ in the frozen protocol; ``level'' means one such $t$.
A statistic ``pooled over the schedule'' sums over levels before taking a root or a
ratio; see the chapter on the statistics of the experiments.

\section{Vectors, matrices, tensors and scalars}

Vectors are bold lowercase ($\va$, $\vx$, $\vz$, $\vect{s}$), matrices and tensors bold
uppercase ($\Sig$, $\Id$, $\Xis$, $\kfour$). Anything indexed down to a scalar loses its
bold: $a_i$ is the $i$-th site of $\va$, $\Sigma_{ij}$ an entry of $\Sig$,
$\kappa_4[ijkl]$ an entry of $\kfour$. The Euclidean norm of a vector is $|\va|$ or
$\lVert\va\rVert$; the Frobenius norm of a matrix or of a symmetric tensor is
$\Frob{\mat{T}}$, with
\begin{equation}
  \Frob{\mat{T}}^2=\sum_{i_1,\dots,i_r}T_{i_1\cdots i_r}^2,
  \label{eq:frob}
\end{equation}
the sum running over \emph{ordered} index tuples. For a symmetric matrix this counts
each off-diagonal entry twice, and for a symmetric 4-tensor it counts an entry with
distinct indices $24$ times. The paper's constants $\tfrac12\Frob{\partial_\theta\Sig_0}^2$
and $\tfrac16\Frob{\kfour}^2$ are stated with this convention, and
\code{src/markov\_score/theory.py} sums the full $(\nsites,\nsites,\nsites,\nsites)$
array; the chapter on moments and cumulants explains how the ordered-index sum turns
into a sum over multisets with weights $4!/\nu!$, which is where the paper's $\tfrac1{24}$
comes from.

$\Id$ is the identity, $\tr$ the trace, $\diag(\vect{v})$ the diagonal matrix with
$\vect{v}$ on the diagonal, $\mat{A}\succ0$ positive definiteness, $\lmax(\mat{A})$ the
largest eigenvalue. $\normal{m}{v}$ is the Gaussian law with mean $m$ and variance (or
covariance matrix) $v$. The symbol $\gamma$ is reserved for the standard Gaussian
density on $\Reals^{\nsites}$, and $\gamma_{\Dt}$ for $\normal{0}{\Dt\Id}$.

\section{Hermite polynomials: probabilists' versus physicists'}

Two families of polynomials carry the name Hermite, and they differ by a rescaling
that changes every constant. The \term{probabilists' Hermite polynomials} are
\begin{equation}
  \Hen{n}(x)=(-1)^n e^{x^2/2}\frac{\dd^n}{\dd x^n}e^{-x^2/2},
  \label{eq:he-def}
\end{equation}
so that $\Hen{0}=1$, $\Hen{1}=x$, $\Hen{2}=x^2-1$, $\Hen{3}=x^3-3x$, $\Hen{4}=x^4-6x^2+3$; they are
monic and orthogonal under the standard Gaussian, with $\E_\gamma[\Hen{n}\Hen{m}]=n!\,\delta_{nm}$.
The \term{physicists' Hermite polynomials} are
$H_n(x)=(-1)^ne^{x^2}\frac{\dd^n}{\dd x^n}e^{-x^2}$, with leading coefficient $2^n$ and
orthogonal under $e^{-x^2}$ with norm $2^nn!\sqrt\pi$. The two are related by
\begin{equation}
  H_n(x)=2^{n/2}\Hen{n}(\sqrt2\,x),\qquad \Hen{n}(x)=2^{-n/2}H_n(x/\sqrt2).
\end{equation}
The paper and this compendium use the probabilists' family throughout, because the
noise is standardised and the orthogonality weight is then exactly the law of the
noise. One notational hazard must be named: the paper writes $H_\nu(\vx)=\prod_i\Hen{\nu_i}(x_i)$
for the \emph{multivariate product} of probabilists' polynomials. That $H_\nu$ is not
a physicists' polynomial, despite the capital letter; it is the product of
probabilists' ones, and $\E_\gamma[H_\nu H_\mu]=\nu!\,\delta_{\nu\mu}$. The chapter on
Gaussian integrals and Hermite polynomials derives all of this.

\section{Standardised white noise and the explicit $\sqrt2$}

The forward process is written, in the paper and in the code
(\code{src/markov\_score/noising.py}), as the Ornstein--Uhlenbeck equation
\begin{equation}
  \dd\vx_t=-\vx_t\,\dd t+\sqrt2\,\dd\vect{W}_t,\qquad \vx_0=\va\sim\Pz,
  \label{eq:ou-sde}
\end{equation}
with $\vect{W}$ a standard Brownian motion, equivalently with a \term{standardised
white noise} $\gwn$ satisfying $\langle\eta_i(t)\eta_j(t')\rangle=\delta_{ij}\delta(t-t')$
and $\dd\vect{W}_t=\gwn(t)\dd t$. The factor $\sqrt2$ is carried explicitly for one
reason: with it, the stationary variance of each coordinate is exactly one. Solving
\cref{eq:ou-sde} by variation of constants,
\begin{equation}
  \vx_t=e^{-t}\va+\sqrt2\int_0^te^{-(t-s)}\dd\vect{W}_s,
  \qquad
  \Var\Bigl(\sqrt2\int_0^te^{-(t-s)}\dd W_s\Bigr)=2\int_0^te^{-2(t-s)}\dd s=1-e^{-2t},
\end{equation}
so that, conditionally on $\va$,
\begin{equation}
  \vx_t\mid\va\ \sim\ \normal{e^{-t}\va}{\Dt\Id},\qquad \Dt=1-e^{-2t},
  \qquad\text{that is}\qquad \vx_t=e^{-t}\va+\sqrt{\Dt}\,\vz,\ \vz\sim\normal{0}{\Id}.
  \label{eq:channel-conv}
\end{equation}
Since $e^{-2t}+\Dt=1$, a coordinate of unit variance keeps unit variance at every $t$,
and $\Pt\to\normal{0}{\Id}$. Two other conventions appear in the literature and are
\emph{not} used here. The physicists' convention absorbs the $\sqrt2$ into the noise
by demanding $\langle\eta(t)\eta(t')\rangle=2\delta(t-t')$; the formulas are the same
but the noise is no longer standardised. The ``variance-preserving'' convention of the
score-based literature writes $\dd\vx=-\tfrac12\beta(t)\vx\,\dd t+\sqrt{\beta(t)}\,\dd\vect W$
with a schedule $\beta(t)$; \cref{eq:ou-sde} is the case $\beta\equiv2$, and the
noise level is then parameterised by $t$ directly. The decay factor is written
$e^{-t}$ in every displayed equation of the paper and is never abbreviated to
$\alpha_t$; inside proofs the paper allows the local abbreviation $\alpha=e^{-t}$,
always introduced in the sentence that uses it, and this compendium does the same.
The \term{signal-to-noise ratio} of the channel is $e^{-2t}/\Dt$, which is
$e^{-2t}(1+O(e^{-2t}))$ at large $t$; this is why the paper's expansions are in
powers of $e^{-2t}$.

\section{What ``the score'' means in the paper}

The word \term{score} is used in two different senses in statistics and in diffusion
modelling, and the paper is careful to use it in one of them. The
\term{diffusion score} is the gradient of the log-density of the \emph{noised data law}
with respect to the \emph{data},
\begin{equation}
  \score(\vx,t)=\nabla_{\vx}\log\Pt(\vx),
\end{equation}
a vector field on $\Reals^{\nsites}$ for each $t$. The \term{statistical score} is the
gradient of a log-likelihood with respect to a \emph{parameter},
$\partial_\theta\log p_\theta(\vx)$, a function of $\theta$ for fixed data. The chapter on
Fisher information and the Cram\'er--Rao bound defines both and keeps them apart; the
paper writes ``the score'' for the diffusion score and ``the likelihood score'' or
``per-sequence scores obtained by Fisher's identity'' for the statistical one.

Within the diffusion sense, four objects have to be distinguished, and the paper's
sentence ``throughout, `the score' means this tree-exact inference, evaluated by
validated quadrature'' names one of them. They are:
\begin{enumerate}[label=(L\arabic*)]
\item exact functional sum-product on the posterior factor graph, which for this model
  is a chain and therefore a tree, so that the forward and backward messages are
  exact and Tweedie's identity $\score(\vx,t)=-(\vx-e^{-t}\pmean{\va})/\Dt$ turns the
  exact posterior mean into the exact score;
\item Gaussian-projected sum-product, in which every message is replaced by the Gaussian
  with its first two moments; the paper's Proposition 1 shows this returns the linear
  MMSE estimator of the covariance twin;
\item the numerical implementation of (L1) on a grid of $\ngrid$ points on
  $[-\halfwidth,\halfwidth]$ with trapezoidal weights, at the frozen setting
  $\ngrid=401$, $\halfwidth=8$, validated by refinement and by an independent brute-force
  enumeration on short chains;
\item the LMMSE estimator of the true covariance, $e^{-t}\Sig_0\Sig_t^{-1}\vx$, and its
  score $\Sgauss=-\Sig_t^{-1}\vx$, which is what (L2) computes.
\end{enumerate}
``The score'' in the paper is (L1) computed as (L3). It is the reference against which
every estimator, whether the EM-fitted kernel or a trained network, is scored; the
estimand of the experiments is the relative error to this reference, pooled over the
schedule. It is never the output of a learned model; those are ``estimates'',
``fits'' or ``arms''. The grid setting is not inherited from the literature but
measured: the chapter on grid quadrature explains what truncation and quadrature
errors are and how the paper controls them, and why a heavy-tailed innovation tests the
domain $\halfwidth$ rather than the spacing $\gridh$.

\section{Other fixed choices}

\begin{itemize}
\item The clean chain has $a_1\sim\normal{0}{1}$, $|\corr|<1$, innovation variance
  $\ivar=1-\corr^2$, so that $\Var(a_i)=1$ and $\Cov(a_i,a_j)=\corr^{|i-j|}$ for every
  innovation law. The frozen configuration is $\corr=0.85$, $\nsites=32$, Laplace
  innovations; every departure from it is stated where it occurs.
\item Innovation families are always \emph{variance-matched}: each is parameterised
  by its variance $\ivar$, and the remaining parameter, if any, is a shape parameter
  (the Student-$t$ degrees of freedom $\nu$, the mixture separation $\kappa$, the
  generalised-Gaussian exponent $\beta$).
\item Cumulants are written $\kappa_m$; the excess kurtosis of a scalar law is
  $\kappa_4/\kappa_2^2$; the fourth-cumulant tensor of a vector is $\kfour$.
\item $\log$ is the natural logarithm; a ``log-slope'' is a derivative or a regression
  coefficient of $\log y$ against $t$, so that $y\propto e^{-8t}$ has log-slope $-8$.
\item British spelling and the serial ``and'' are used; no number is typed from a
  measurement.
\end{itemize}

\where{Every chapter. The symbol file is \code{paper/notation.tex}, read verbatim by
\code{paper/main.tex} and by \code{reports/compendium/compendium.tex}; the channel
convention is implemented in \code{src/markov\_score/noising.py} (the functions
\code{alpha\_delta} and \code{ou\_noise\_sample}); the Hermite convention in
\code{src/markov\_score/theory.py} (the function \code{he}); the frozen configuration
in \code{experiments/frozen\_config.py} and the comparison protocol in
\code{experiments/protocol.py}. The paper's Section 2 fixes the chain and the channel,
its Section 3 the Hermite notation, and its Appendix B the grid.}

\paragraph{References.}
Y. Song, J. Sohl-Dickstein, D. P. Kingma, A. Kumar, S. Ermon and B. Poole (2021),
``Score-based generative modeling through stochastic differential equations'', ICLR.
M. M\'ezard and A. Montanari (2009), \emph{Information, Physics, and Computation},
Oxford University Press.
G. Szeg\H{o} (1975), \emph{Orthogonal Polynomials}, 4th edition, American Mathematical
Society (Chapter V for the two Hermite conventions).
C. W. Gardiner (2009), \emph{Stochastic Methods}, 4th edition, Springer (Chapter 4 for
the Ornstein--Uhlenbeck process and the white-noise conventions).
B. Efron (2011), ``Tweedie's formula and selection bias'', Journal of the American
Statistical Association 106, 1602--1614.
